\section{练习}

\begin{question}\begin{drill}[电子绕核运动的等效电流]\label{c1:drill:1}
电子以速率 $v$ 沿半径 $R$ 的圆轨道运动。求通过轨道上任一固定截面的平均等效电流及方向。
\end{drill}
\typeout{ECFIG:fig:c1:drill:1:s}\fig{fig:c1:drill:1:s}{题设：电子绕核运动的等效电流}{\begin{tikzpicture}[scale=.88]
\draw[world](0,0)circle(1.3);\draw[flow](1.3,0)arc(0:65:1.3);\node at(1.8,.8){$v$};\draw[dim](0,0)--(-1.3,0)node[midway,above]{$R$};\node at(0,0){$+$};
\end{tikzpicture}}
\end{question}
\answerlink{c1:drill:1}
\begin{question}\begin{drill}[固定碰撞间隔与平均漂移]\label{c1:drill:2}
某简化碰撞模型规定电子两次碰撞严格相隔 $T$，碰撞后定向速度为零，热运动均值为零。电子质量为 $m$、电荷为 $-e$、数密度为 $n$，在均匀电场中运动。求平均漂移速度和电阻率，解释其中的因子 $2$。
\end{drill}
\typeout{ECFIG:fig:c1:drill:2:s}\fig{fig:c1:drill:2:s}{题设：固定碰撞间隔与平均漂移}{\begin{tikzpicture}[scale=.88]
\draw[axis](-.2,0)--(4.8,0)node[right]{$t$};\draw[axis](0,-.2)--(0,2.6)node[above]{$|v_{\rm d}|$};\draw[world](0,0)--(1.3,1.8)--(1.3,0)--(2.6,1.8)--(2.6,0)--(3.9,1.8);\node[below]at(1.3,0){$T$};\node at(2.6,2.4){电子定向速率示意};
\end{tikzpicture}}
\end{question}
\answerlink{c1:drill:2}
\begin{question}\begin{drill}[圆台形细导体的电阻]\label{c1:drill:3}
均匀材料电阻率为 $\rho$。长 $l$ 的细圆台两端半径为 $a,b$，半径沿轴线性变化，$|b-a|\ll l$。忽略端部效应，求轴向电阻。
\end{drill}
\typeout{ECFIG:fig:c1:drill:3:s}\fig{fig:c1:drill:3:s}{题设：圆台形细导体的电阻}{\begin{tikzpicture}[scale=.88]
\draw[world](0,-.4)--(4,-.9)--(4,.9)--(0,.4)--cycle;\draw[guide](0,0)--(4,0);\draw[dim](0,-1.4)--(4,-1.4)node[midway,below]{$l$};\node[left]at(0,0){$a$};\node[right]at(4,0){$b$};
\end{tikzpicture}}
\end{question}
\answerlink{c1:drill:3}
\begin{question}\begin{drill}[同心球电极之间的电阻]\label{c1:drill:4}
两同心金属球面半径为 $a,b$，中间充满电阻率 $\rho$ 的材料。求两球面间电阻及 $b\to\infty$ 的极限。
\end{drill}
\typeout{ECFIG:fig:c1:drill:4:s}\fig{fig:c1:drill:4:s}{题设：同心球电极之间的电阻}{\begin{tikzpicture}[scale=.88]
\draw[world](0,0)circle(.65);\draw[world](0,0)circle(1.6);\draw[dim](0,0)--(.65,0)node[midway,above]{$a$};\draw[dim](0,0)--(0,-1.6)node[midway,left]{$b$};\node at(3,0){球对称};
\end{tikzpicture}}
\end{question}
\answerlink{c1:drill:4}
\begin{question}\begin{drill}[两个半球电极的近似总电阻]\label{c1:drill:5}
均匀大地电阻率 $\rho$，两相同半球电极半径 $a$，球心在地表且相距 $d\gg a$。电流从 $P$ 注入，从 $Q$ 收集，求忽略电极表面非均匀修正时的两电极间电阻。
\end{drill}
\typeout{ECFIG:fig:c1:drill:5:s}\fig{fig:c1:drill:5:s}{题设：两个半球电极的近似总电阻}{\begin{tikzpicture}[scale=.88]
\draw[world](-1,0)--(5,0);\draw[world](-.4,0)arc(180:360:.4);\draw[world](3.6,0)arc(180:360:.4);\node[above]at(0,0){$P$};\node[above]at(4,0){$Q$};\draw[dim](0,.8)--(4,.8)node[midway,above]{$d$};\node at(2,-1.3){半径均为 $a$};
\end{tikzpicture}}
\end{question}
\answerlink{c1:drill:5}
\begin{question}\begin{drill}[电阻丝拉长后的导电变化]\label{c1:drill:6}
一根电阻丝长度 $l$、面积 $S$、电阻 $R$。在载流子总数、体积与电阻率保持不变的理想条件下均匀拉长为 $kl$。若仍接同一理想电压源，求电阻、电流、漂移速率和总发热功率的变化。
\end{drill}
\typeout{ECFIG:fig:c1:drill:6:s}\fig{fig:c1:drill:6:s}{题设：电阻丝拉长后的导电变化}{\begin{tikzpicture}[scale=.88]
\draw[world](0,0)rectangle(4,1.1);\draw[guide](1.6,0)--(1.6,1.1);\draw[flow](.3,.55)--(1.3,.55)node[midway,above]{$I$};\draw[dim](0,-.5)--(4,-.5)node[midway,below]{$l$};\node[left]at(0,.55){$A$};\node[right]at(4,.55){$B$};\node[above]at(2,1.1){$S,\rho$};\node at(2,2){拉长至 $kl$，体积不变};
\end{tikzpicture}}
\end{question}
\answerlink{c1:drill:6}
\begin{question}\begin{drill}[电池充电时的功率]\label{c1:drill:7}
电池电动势 $\mathcal E=12\,\mathrm V$，内阻 $r=0.50\,\Omega$。外装置使 $I=2.0\,\mathrm A$ 的电流从电池正端流入。求端电压、输入功率、内部发热功率及理想化学储能速率。
\end{drill}
\typeout{ECFIG:fig:c1:drill:7:s}\fig{fig:c1:drill:7:s}{题设：电池充电时的功率}{\begin{tikzpicture}[scale=.88]
\draw(0,0)to[battery1,invert,l={$\mathcal E$}](0,2.4)to[R,l={$r$}](3,2.4);\draw[flow](2,2.9)--(.8,2.9)node[midway,above]{$I$};\node[above left]at(0,2.4){$+$};\node[below left]at(0,0){$-$};\node[right]at(3,2.4){$A$};\draw(0,0)--(3,0)node[right]{$B$};
\end{tikzpicture}}
\end{question}
\answerlink{c1:drill:7}
\begin{question}\begin{drill}[恒流加热的稳态与稳定性]\label{c1:drill:8}
导体电阻 $R(T)=R_0[1+\alpha(T-T_0)]$，$\alpha>0$。环境温度 $T_0$，散热功率 $k(T-T_0)$，热容量 $C$ 恒定。接恒流源 $I$。求可能的稳定温升，并说明何时线性模型没有正的稳态解。
\end{drill}
\typeout{ECFIG:fig:c1:drill:8:s}\fig{fig:c1:drill:8:s}{题设：恒流加热的稳态与稳定性}{\begin{tikzpicture}[scale=.88]
\node[draw,minimum width=2cm,minimum height=1cm]at(1,0){$R(T),C$};\draw[flow](-1,0)--(0,0)node[midway,above]{$I$};\draw[flow](2,0)--(4,0)node[midway,above]{散热};\node at(3,-.65){$k(T-T_0)$};
\end{tikzpicture}}
\end{question}
\answerlink{c1:drill:8}
\begin{question}\begin{drill}[质子束中的线密度变化]\label{c1:drill:9}
采用非相对论模型，忽略束内相互作用。质子由近似静止开始在均匀场中加速，束电流恒为 $1.0\,\mathrm{mA}$。求每秒打靶粒子数；在离源 $L$ 与 $4L$ 处取相同的很短空间段，求段内粒子数比。质子电荷 $e=1.60\times10^{-19}\,\mathrm C$。
\end{drill}
\typeout{ECFIG:fig:c1:drill:9:s}\fig{fig:c1:drill:9:s}{题设：质子束中的线密度变化}{\begin{tikzpicture}[scale=.88]
\draw[flow](0,0)--(6,0);\draw[world](1,-.6)rectangle(1.3,.6);\draw[world](4,-.6)rectangle(4.3,.6);\node at(1,-1){$L$};\node at(4,-1){$4L$};\node at(3,1){均匀加速，束流 $I$ 恒定};
\end{tikzpicture}}
\end{question}
\answerlink{c1:drill:9}
\begin{question}\begin{drill}[持续电流给出的电阻率上限]\label{c1:drill:10}
闭合细金属环初电流 $I_0=100\,\mathrm A$，一年内未观察到超过 $\Delta I=1.0\,\mathrm{mA}$ 的衰减。采用无自感反馈的有效阻力模型，$n=8.00\times10^{20}\,\mathrm{m^{-3}}$，求此模型下电阻率上限。再说明真实超导环用此上限时应注意什么。
\end{drill}
\typeout{ECFIG:fig:c1:drill:10:s}\fig{fig:c1:drill:10:s}{题设：持续电流给出的电阻率上限}{\begin{tikzpicture}[scale=.88]
\draw[world](0,0)circle(1.6);\draw[flow](1.6,0)arc(0:75:1.6);\node at(0,0){$I(t)$};\node at(3,0){一年内 $\Delta I\le1\,\mathrm{mA}$};
\end{tikzpicture}}
\end{question}
\answerlink{c1:drill:10}
